% Источники: exam/materials/преподаватель/2020/27-04-2020-Model_L10_3_.pdf; exam/materials/моделирование.pdf, раздел 35; exam/materials/прошлогодние ответы/Ответы_Моделирование.pdf, раздел 35.
\section{Ранговая корреляция Кенделла.}

\subsection*{Коэффициент ранговой корреляции Кенделла}

\[
  \tau = \frac{4P}{n(n-1)} - 1,
\]
где $P$ — число \textbf{совпадений рангов}, то есть число конкордантных пар.

Связь с числом \textbf{конкордантных пар} $C$ и дискордантных $D$:
\[
  P = C, \qquad
  \tau = \frac{C - D}{\dfrac{n(n-1)}{2}}
       = 1 - \frac{4D}{n(n-1)}.
\]

\subsection*{Алгоритм}

\begin{enumerate}
  \item Упорядочить наблюдения по возрастанию ранга $R_i$ (по $X$).
  \item Для каждого $i$ подсчитать число $j > i$, для которых $S_j > S_i$ (конкордантные, совпадение рангов) и $S_j < S_i$ (дискордантные).
  \item $C$ — число конкордантных пар, $D$ — число дискордантных.
  \item Подставить $P = C$ в лекционную формулу.
\end{enumerate}

\subsection*{Интерпретация}

$\tau \in [-1, 1]$: значение $+1$ — полная прямая монотонная зависимость, $-1$ — обратная.

\subsection*{Пример}

\begin{center}
\begin{tabular}{cccc}
  \hline
  $i$ & $X_i$ & $Y_i$ & $R_i$ \\
  \hline
  1 & 1 & 3 & 1 \\
  2 & 2 & 1 & 2 \\
  3 & 3 & 4 & 3 \\
  4 & 4 & 2 & 4 \\
  \hline
\end{tabular}
\end{center}

Ранги $Y$ в порядке $R_i$: последовательность $S = (3, 1, 4, 2)$.

Подсчёт дискордантных пар:
\begin{itemize}
  \item $S_1 = 3$: ниже стоят $S_2=1$ ($<3$, $D{+}1$), $S_3=4$ ($>3$, $C{+}1$), $S_4=2$ ($<3$, $D{+}1$).
  \item $S_2 = 1$: ниже стоят $S_3=4$ ($>1$, $C{+}1$), $S_4=2$ ($>1$, $C{+}1$).
  \item $S_3 = 4$: ниже стоит $S_4=2$ ($<4$, $D{+}1$).
\end{itemize}

$C = 3$, $D = 3$, $P=C=3$, $n = 4$.
\[
  \tau = \frac{4P}{n(n-1)} - 1
       = \frac{4\cdot 3}{4\cdot 3} - 1 = 0.
\]

Вывод: монотонная зависимость между $X$ и $Y$ не обнаружена.

\clearpage
